\documentclass[12pt, twoside]{article}
\input{../preamble}

\fancyhead[LE]{\thepage}
\fancyhead[RO]{Markscheme}
\fancyhead[LO]{La Scuola d'Italia / Dr.\ Huson / IB Math (c) \\* 14 March 2026}

\begin{document}
\subsubsection*{7.11 PreQuiz: Logarithmic and Exponential Functions ---
Markscheme}

\begin{enumerate}[itemsep=0.3cm]

%--- Problem 1: Quarterly compound interest and continuous equivalent [6 marks] ---
\item[1.] [Maximum mark: 6]

Noah invests \$3200 at $7.2\,\%$ per annum, compounded quarterly.

\medskip

\textbf{(a)} Quarterly rate $=\dfrac{0.072}{4}=0.018$. \hfill \textbf{M1}

$V(t)=3200(1.018)^{4t}$ \hfill \textbf{A1}

\medskip

\textbf{(b)} Since $(1.018)^{4t}=e^{4t\ln(1.018)}$,
\hfill \textbf{M1}

$V(t)=3200e^{kt}$ where $k=4\ln(1.018)$. \hfill \textbf{A1}

\medskip

\textbf{(c)} Need the smallest integer $n$ with
$3200(1.018)^{4n}>5000$. \hfill \textbf{M1}

$n>6.25\ldots$, so the smallest full-year value is $n=7$ years.
\hfill \textbf{A1}

\emph{Alternative method (table): show both adjacent rows
$3200(1.018)^{24}=4910.48$ (at $n=6$) and $3200(1.018)^{28}=5273.18$
(at $n=7$), then state $n=7$. A single-row table earns full marks but
should receive a feedback note recommending the adjacent row to
justify the crossover.}

\newpage

%--- Problem 2: Continuous exponential growth [6 marks] ---
\item[2.] [Maximum mark: 6]

$U(t)=9000e^{kt}$ and $U(3)=18\,500$.

\medskip

\textbf{(a)} $9000e^{3k}=18\,500$, so $e^{3k}=\dfrac{18500}{9000}$.
\hfill \textbf{M1}

$k=\dfrac{1}{3}\ln\!\left(\dfrac{18500}{9000}\right)
=0.2401\ldots$ \hfill \textbf{A1}

\medskip

\textbf{(b)} $U(7)=9000e^{7k}$, using their value of $k$.
\hfill \textbf{M1}

$U(7)=48\,351.58\ldots$, so $48\,352$ users. \hfill \textbf{A1}

\emph{Accept $48\,290$ (obtained by substituting rounded $k=0.240$)
with a feedback note: do not round intermediate values; carry full
precision through and round only at the end.}

\medskip

\textbf{(c)} Set $9000e^{kt}=50\,000$. \hfill \textbf{M1}

$t=\dfrac{\ln(50000/9000)}{k}=7.14$ years. \hfill \textbf{A1}

\emph{Accept $7.15$ (obtained by substituting rounded $k=0.240$) with
the same feedback note about not rounding intermediate values.}

\end{enumerate}
\end{document}
